\section{Scope Limits and Non-Claims}

This work is strictly limited in scope. The following non-claims are stated
explicitly to prevent misinterpretation of the results.

\paragraph{No Prescriptive or Normative Claims.}
The results presented here do not prescribe actions, strategies, governance
mechanisms, or optimization criteria. No statements of the form ``should,''
``ought,'' or ``better'' are made or implied. The analysis is descriptive and
structural only.

\paragraph{No Recovery or Turnaround Claims.}
This work does not address recovery, reconstitution, or restoration of
enterprise identity. In particular, no mechanisms for reversing identity loss
are proposed, and no claims are made regarding the feasibility of turnaround
without introducing new admissible structure, which is explicitly excluded by
assumption.

\paragraph{No Managerial, Cultural, or Narrative Semantics.}
Enterprise identity is not interpreted in psychological, cultural, branding, or
managerial terms. The formal notion of identity used here is a structural
property of the continuum $K_{\mathrm{EA}}$ and must not be conflated with
self-descriptions, stakeholder narratives, or symbolic representations.

\paragraph{No Extension of OC or ETS.}
The analysis introduces no new primitives, operators, thresholds, or phase
classes beyond those already present in the Ontology of Continua and the
executable specification $\mathrm{ETS}.K_{\mathrm{EA}}$. All results hold only
under the closed dynamics of these frameworks.

\paragraph{Domain of Validity.}
The conclusions apply exclusively to enterprises that admit a representation as
a continuum $K_{\mathrm{EA}}$ with non-trivial cycle structure. Enterprises or
organizations that cannot be so modeled lie outside the domain of this work.

These scope limits are essential to prevent misinterpretation of identity loss as
a managerial failure, cultural deficit, or optimization problem. The phenomenon
analyzed here is a structural phase transition, and nothing more.
